\documentclass{article}

%===============================
%
%          📦 Paquetes
%
%===============================

% === Configuración del documento y tipografía ===
\usepackage[a4paper, top=3cm, bottom=2.5cm, left=2.5cm, right=2.5cm]{geometry}
\usepackage[spanish]{babel}
\usepackage[utf8]{inputenc}
\usepackage{newtxtext} % Times New Roman-like font
%\usepackage{newtxmath} % Matching math font

% === Gráficos y representaciones visuales ===
\usepackage{tikz}
\usepackage{pgfplots}
\pgfplotsset{compat=1.18}
\usepackage{graphicx}
\usepackage{subcaption}
\usepackage{booktabs}
\pgfplotsset{compat=1.18}

% === Estilo y personalización ===
\usepackage{titling}
\usepackage{titlesec}
\usepackage{tocloft}
\usepackage{fancyhdr}
\usepackage{setspace}
\usepackage{caption}

% === Matemáticas y símbolos ===
\usepackage{amsmath}
\usepackage{amssymb}
\usepackage{cancel}

% === Estructura y disposición del contenido ===
\usepackage{array}
\usepackage{multicol}
\usepackage{float}
\usepackage{indentfirst}
\usepackage{titlesec}

% === Hipervínculos y referencias ===
\usepackage{tocbibind}
\usepackage{hyperref}

%===============================
%
%          ⬆️⬇️ Cabeceras
%
%===============================

\pagestyle{fancy}
\fancyhf{}
\fancyhead[L]{\includegraphics[width=2cm]{assets/logo.png}}
\fancyhead[R]{\textit{Mecánica Clásica - 44280}}
\fancyfoot[R]{\thepage}
\setlength{\headheight}{22.9188pt}
\addtolength{\topmargin}{-5.64944pt}

%======================================
%
%          📑 Espaciado de Párrafos
%
%======================================

\setlength{\parskip}{1em}
\setlength{\parindent}{0pt}
\titlespacing*{\section}
  {0pt}{1.5ex plus 1ex minus .2ex}{0.1ex plus .1ex}

\titlespacing*{\subsection}
  {0pt}{1.2ex plus 0.8ex minus .2ex}{0.2ex plus .1ex}

\titlespacing*{\subsubsection}
  {0pt}{1ex plus 0.5ex minus .2ex}{0.1ex plus .1ex}



%======================================
%
%          📕 Estrucura del título
%
%======================================

% \title{
%   \textbf{Relación entre la capacitación continua de los trabajadores y su productividad en una empresa de tecnología en Perú durante el año 2025}
% }
% \author{
%   \begin{tabular}{ll}
%     \textbf{Huatay Salcedo, Luis} & \texttt{U24218809}
%   \end{tabular}
% }

\begin{document}

%======================================
%
%          📚 Inicio del documento
%
%======================================

% Introducción
\vspace*{\fill}

\section{Introducción}

  El presente proyecto tiene como finalidad aplicar la Metodología de Gestión de Proyectos Clásica, siguiendo el marco de trabajo propuesto por el Project Management Institute (PMI) y utilizando herramientas del Framework del PMBOK. Durante la primera parte de la unidad, se desarrollará un marco predictivo para la gestión de un proyecto empresarial, abarcando las fases de Iniciación y Planificación.

  Para ello, se seleccionará una empresa y área específica, describiendo su misión, visión y objetivos estratégicos, así como el alineamiento del proyecto con dichos objetivos. Posteriormente, se documentarán los procesos de iniciación (objetivos, alcance inicial, gestión de interesados y acta de constitución) y planificación (gestión de integración, alcance, cronograma, costes, calidad, recursos, comunicaciones, riesgos y adquisiciones).

  Finalmente, se presentará un informe detallado que incluya todos los elementos mencionados, demostrando la capacidad de aplicar la metodología de gestión de proyectos en un entorno empresarial real.

\vspace*{\fill}

% Fundamentos Físicos
\section{Fundamentos Físicos}

Los fundamentos físicos del proyecto se centran los que intervienen en el diseño y funcionamiento del mecanismo \textit{Theo Jansen}, estos fundamentos incluyen principios de mecánica clásica como cinemática y dinámica, que son esenciales para comprender cómo el mecanismo puede aprovechar la energía cinética del viento para moverse y con esto, generar números aleatorios. A continuación, se describen los conceptos clave que sustentan el proyecto:
\begin{itemize}
  \item \textbf{Cinemática:} La cinemática es la rama de la mecánica que estudia el movimiento de los cuerpos sin considerar las fuerzas que lo producen. En el caso del mecanismo \textit{Theo Jansen}, la cinemática se aplica para analizar el movimiento rotacional que se traduce en traslacional. Se utilizan conceptos como posición, velocidad y aceleración para describir el movimiento del mecanismo en función del tiempo.

  Algunas de las expresiones matemáticas utilizadas en la cinemática incluyen:

\textbf{Ecuación de la posición:}
\begin{equation}
x(t) = x_0 + v_0 t + \frac{1}{2} a t^2
\end{equation}
Donde: 

\begin{tabular}{@{}ll}
  \(x(t)\): & Posición en el tiempo \(t\) \\
  \(x_0\):  & Posición inicial \\
  \(v_0\):  & Velocidad inicial \\
  \(a\):    & Aceleración \\
\end{tabular}

\textbf{Razón de cambio de la posición o velocidad:}
\begin{equation}
v(t) = v_0 + a t
\end{equation}
Donde:

\begin{tabular}{@{}ll}
  \(v(t)\): & Velocidad en el tiempo \(t\) \\
  \(v_0\):  & Velocidad inicial \\
  \(a\):    & Aceleración \\
\end{tabular}

\textbf{Razón de cambio de la velocidad o aceleración:}
\begin{equation}
a(t) = \frac{dv}{dt}
\end{equation}

Donde:

\begin{tabular}{@{}ll}
  \(a(t)\): & Aceleración en el tiempo \(t\) \\
\end{tabular}


  \item \textbf{Dinámica:} La dinámica es la rama de la mecánica que estudia las fuerzas y sus efectos sobre el movimiento de los cuerpos. En este proyecto, la dinámica se aplica para comprender cómo las fuerzas del viento afectan el movimiento del mecanismo y cómo estas fuerzas pueden ser aprovechadas para generar números aleatorios. Se utilizan conceptos como fuerza, masa y aceleración para analizar el comportamiento del mecanismo bajo la influencia del viento.
  
  La expresión fundamental de la dinámica es la segunda ley de Newton, que establece la relación entre la fuerza neta aplicada a un cuerpo, su masa y su aceleración:

\textbf{Segunda ley de Newton:}
\begin{equation}
\vec{F} = m\dfrac{\Delta \vec{V}}{\Delta t} = m\vec{a}
\end{equation}
Donde:

\begin{tabular}{@{}ll}
  \(\vec{F}\): & Fuerza neta aplicada al cuerpo \\
  \(m\):      & Masa del cuerpo \\
  \(\Delta \vec{V}\): & Cambio en la velocidad del cuerpo \\
  \(\Delta t\): & Intervalo de tiempo durante el cual se aplica la fuerza \\
  \(\vec{a}\): & Aceleración del cuerpo \\
\end{tabular}


  \item \textbf{Energía cinética:} La energía cinética es la energía asociada al movimiento de un cuerpo. En el mecanismo \textit{Theo Jansen}, la energía cinética del viento se convierte en movimiento mecánico, lo que permite que el mecanismo se desplace y genere números aleatorios. La energía cinética se calcula utilizando la fórmula:
  
\textbf{Energía cinética:}
\begin{equation}
E_k = \frac{1}{2} m v^2
\end{equation}
Donde:

\begin{tabular}{@{}ll}
  \(E_k\): & Energía cinética del cuerpo \\
  \(m\):   & Masa del cuerpo \\
  \(v\):   & Velocidad del cuerpo \\
\end{tabular}

  \item \textbf{Entropía:} La entropía es una medida del desorden o aleatoriedad en un sistema. En el contexto de este proyecto, la entropía se refiere a la variabilidad en el movimiento del mecanismo lo que causará intervalos de movimiento y detención, lo que a su vez generará números aleatorios. La entropía se relaciona con la cantidad de información que se puede obtener de un sistema y es fundamental para garantizar la aleatoriedad de los números generados.
  
\textbf{Entropía:}
\begin{equation}
S = -\sum_{i=1}^{n} p_i \log(p_i)
\end{equation}
Donde:  

\begin{tabular}{@{}ll}
  \(S\): & Entropía del sistema \\
  \(p_i\): & Probabilidad de ocurrencia del evento \(i\) \\
  \(n\): & Número total de eventos posibles \\
\end{tabular}

  \item \textbf{Movimiento circular uniforme:} El movimiento circular uniforme es un tipo de movimiento en el que un objeto se desplaza a lo largo de una trayectoria circular con una velocidad constante. En el caso del mecanismo \textit{Theo Jansen}, el movimiento circular uniforme se utiliza para generar el movimiento rotacional necesario para el funcionamiento del mecanismo. Este tipo de movimiento se caracteriza por una velocidad angular constante y una aceleración centrípeta que mantiene al objeto en su trayectoria circular.
  
  Así mismo, para nuestro mecamismo, este utilizará una ruleta de números aleatorios que se detendrá en un número al azar cada vez que el mecanismo se detenga. La ruleta estará dividida en segmentos numerados, y la probabilidad de que el mecanismo se detenga en un número específico dependerá de la distribución de los segmentos y de la entropía del viento.

  Algunas de las expresiones matemáticas utilizadas en el movimiento circular uniforme incluyen:

  \begin{itemize}
    \item \textbf{Velocidad angular:}
  \begin{equation}
  \omega = \frac{\Delta \theta}{\Delta t}
  \end{equation}
  Donde:
  
  \begin{tabular}{@{}ll}
    \(\omega\): & Velocidad angular del objeto \\
    \(\Delta \theta\): & Cambio en el ángulo de posición del objeto \\
    \(\Delta t\): & Intervalo de tiempo durante el cual ocurre el cambio \\
  \end{tabular}

    \item \textbf{Velocidad tangencial:}
  \begin{equation}
  v = r \omega
  \end{equation}
  Donde:

  \begin{tabular}{@{}ll}
    \(v\): & Velocidad tangencial del objeto \\
    \(r\): & Radio de la trayectoria circular \\
    \(\omega\): & Velocidad angular del objeto \\
  \end{tabular}

    \item \textbf{Frecuencia:}
  \begin{equation}
  f = \frac{1}{T}
  \end{equation}
  Donde:

  \begin{tabular}{@{}ll}
    \(f\): & Frecuencia del movimiento circular \\
    \(T\): & Período del movimiento circular (tiempo para completar una vuelta) \\
  \end{tabular}
  
  \end{itemize}
\end{itemize}

% Diseño del mecanismo
\section{Metodología}

El proyecto se lleva a cabo siguiendo varias etapas clave que cincluyen la investigación, el diseño, la construcción y la prueba del mecanismo. Cada etapa se desarrolla de manera sistemática para garantizar que se cumplan los objetivos del proyecto y se obtengan resultados significativos. 

Concluida la etapa de investigación, a continuación se detallan las siguientes etapas de la metodología empleada:

\subsection{Diseño del mecanismo}

El diseño del mecanismo \textit{Theo Jansen} se basa en los principios de la mecánica clásica y la cinética. Las partes importantes del mecanismo incluyen:

\begin{itemize}
  \item \textbf{Ejes y engranajes:} Estos componentes son fundamentales para transmitir el movimiento del viento al mecanismo. Los ejes permiten que las partes móviles giren, mientras que los engranajes ayudan a convertir el movimiento rotacional en movimiento lineal, lo que es esencial para el funcionamiento del mecanismo.
  \item \textbf{Estructura de soporte:} La estructura del mecanismo debe ser lo suficientemente robusta para soportar las fuerzas generadas por el viento y el movimiento del mecanismo y al mismo tiempo ser lo suficientemente ligera para permitir que el viento la mueva. Esto se logra mediante el uso de materiales adecuados y un diseño estructural eficiente.
  \item \textbf{Articulaciones:} Se refiere a las conexiones entre las diferentes partes del mecanismo que permiten el movimiento relativo entre ellas. Estas articulaciones son cruciales para garantizar que el mecanismo funcione de manera fluida y eficiente.
\end{itemize}

Una de las partes fundamentales, luego de diseñar estos componentes, son las patas del mecanismo, que permiten que se mueva y camine de manera eficiente y fluida, según el creador del mecanismo, Theo Jansen. Este simula el movimiento de la pata de un animal y ha sido perfeccionado durante los últimos 10 años mediante un algoritmo evolutivo. Según \textit{Bustamante T. (2016)}, el criterio principal para el desarrollo de estas patas es el rendimiento de los elementos en la tarea encomendada, utilizando los errores y mejoras de las evoluciones para optimizar el diseño en cada iteración. \cite{tellez2016diseno}

Este algoritmo evolutivo entrega al final de sus múltiples iteraciones las medidas de las patas del mecanismo, lo que su creador denomina \textbf{\textit{Los números sagrados}}

\begin{table}[H]
  \centering
  \caption{Valores asignados a cada letra}
  \begin{tabular}{cc|cc}
    \toprule
    \textbf{Letra} & \textbf{Valor} & \textbf{Letra} & \textbf{Valor} \\
    \midrule
    a & 38.0 & g & 36.7 \\
    b & 41.5 & h & 65.7 \\
    c & 39.3 & i & 49.0 \\
    d & 40.1 & j & 50.0 \\
    e & 55.8 & k & 61.9 \\
    f & 39.4 & l & 7.8 \\
    \multicolumn{4}{c}{\textbf{m = 15.0}} \\
    \bottomrule
  \end{tabular}
\end{table}

Estos valores hacen posible un movimiento fluido y áltamente eficiente del mecanismo, permitiendo que se mueva con la energía cinética del viento.

De acuerdo a esta lista es necesario poder escalar las medidas de las patas del mecanismo para que se ajusten a las dimensiones del proyecto. Para ello, se establece un factor de escala que es la relación entre las dimensiones del mecanismo y las dimensiones reales de las patas.

Como medida base se toma la longitud más grande \textbf{h} que pretende ser de 65.7 cm pero se escalará a 21.5 cm de acuerdo a la siguiente expresión:
\begin{equation}
  \text{Factor de escala} = \frac{\text{Longitud deseada}}{\text{Longitud original}} = \frac{21.5}{65.7} \approx 0.33
\end{equation}

Es entonces de acuerdo a este factor de escala que se multiplicarán las medidas de las patas del mecanismo para obtener las dimensiones reales de las patas del mecanismo obteniendo así las longitudes entre ejes de rotación de las articulaciones que se usarán para el proyecto. A continuación se muestran las medidas escaladas de las patas del mecanismo \textit{Theo Jansen} y al costado el margen de medio centímetro en cada extremo que asegura espacio para las uniones de las patas con los ejes de rotación:

\begin{multicols}{2}
  \begin{table}[H]
    \centering
    \caption{Longitud entre ejes de rotación}
    \begin{tabular}{cc|cc}
      \toprule
      \textbf{Letra} & \textbf{Valor} & \textbf{Letra} & \textbf{Valor} \\
      \midrule
      a & 12.5 & g & 12.1 \\
      b & 13.7 & h & 21.7 \\
      c & 13.0 & i & 16.2 \\
      d & 13.2 & j & 16.5 \\
      e & 18.4 & k & 20.4 \\
      f & 13.0 & l & 2.6 \\
      \multicolumn{4}{c}{\textbf{m = 5}} \\
      \bottomrule
    \end{tabular}
  \end{table}
  
  \begin{table}[H]
    \centering
    \caption{Longitud completa con margen}
    \begin{tabular}{cc|cc}
      \toprule
      \textbf{Letra} & \textbf{Valor} & \textbf{Letra} & \textbf{Valor} \\
      \midrule
      a & 13.5 & g & 13.1 \\
      b & 14.7 & h & 22.7 \\
      c & 14.0 & i & 17.2 \\
      d & 14.2 & j & 17.5 \\
      e & 19.4 & k & 21.4 \\
      f & 14.0 & l & 3.6 \\
      \multicolumn{4}{c}{\textbf{m = 6}} \\
      \bottomrule
    \end{tabular}
  \end{table}
\end{multicols}

Una vez encontradas estas medidas se pretenden ensamblar de la siguiente manera:
\begin{figure}[H]
  \centering
  \caption{Diagrama de las patas del mecanismo \textit{Theo Jansen} con los números sagrados} 
  \label{fig:diagrama_barras}
  \begin{tikzpicture}[scale=0.1, thick, every node/.style={scale=1}]

  \coordinate (P1) at (22,30);
  \coordinate (P2) at (55,30);
  \coordinate (P3) at (50,38);
  \coordinate (P4) at (50,50);
  \coordinate (P5) at (15,55);
  \coordinate (P6) at (-20,20);
  \coordinate (P7) at (-10,-20);
  \coordinate (P8) at (30,-10);
  \coordinate (P9) at (40,-60);

  \draw (P1) -- (P2) node[midway,above] {\small a};
  \draw (P2) -- (P3) node[midway,left] {\small l};
  \draw (P3) -- (P4) node[midway,right] {\small m};
  \draw (P4) -- (P5) node[midway,above] {\small j};
  \draw (P5) -- (P1) node[midway,left] {\small b};
  \draw (P5) -- (P6) node[midway,right] {\small e};
  \draw (P6) -- (P1) node[midway,above] {\small d};
  \draw (P6) -- (P7) node[midway,right] {\small f};
  \draw (P1) -- (P8) node[midway,above left] {\small c};
  \draw (P7) -- (P8) node[midway,above] {\small g};
  \draw (P7) -- (P9) node[midway,above] {\small h};
  \draw (P9) -- (P8) node[midway,left] {\small i};
  \draw (P8) -- (P4) node[midway,left] {\small k};
  
  \filldraw (P1) circle (23pt) node[above right] {\small 1};
  \filldraw (P2) circle (23pt) node[above right] {\small 2};
  \filldraw (P3) circle (23pt) node[above right] {\small 3};
  \filldraw (P4) circle (23pt) node[above right] {\small 4};
  \filldraw (P5) circle (23pt) node[above right] {\small 5};
  \filldraw (P6) circle (23pt) node[above left] {\small 6};
  \filldraw (P7) circle (23pt) node[above right] {\small 7};
  \filldraw (P8) circle (23pt) node[above right] {\small 8};
  \filldraw (P9) circle (23pt) node[above right] {\small 9};
  \end{tikzpicture}
\end{figure}


\subsection{Materiales}

Dentro del marco del proyecto, se emplean diversos materiales para la construcción del mecanismo \textit{Theo Jansen}. Estos materiales son seleccionados fueron esocgidos por su durabilidad, flexibilidad y ligeresa, lo que permite que el mecanismo funcione de manera eficiente y efectiva. Los materiales utilizados incluyen:

\begin{itemize}
  \item \textbf{Sorbetes de plástico:} Estos sorbetes son utilizados para las patas del mecanismo, ya que son ligeros y flexibles, lo que permite que el mecanismo se mueva de manera eficiente. Además, los sorbetes son fáciles de conseguir y económicos, lo que los convierte en una opción ideal. Además de ser increíblemente resistentes ante la presión y torsión.
  
  \item \textbf{Hilo de coser o pavilo:} Este hilo se utiliza para unir las diferentes partes de la pata, hacen posible rotar a la articulación de manera que no es afectada por la fricción.
  
  \item \textbf{Palillos de manera:} Es importante fijar las patas a un eje central, para lo cual se utilizan palillos de madera. Estos palillos son resistentes y permiten unir las patas mediante el cigüeñal del mecanismo que convertirá el movimiento rotacional en movimiento lineal.1
  
\end{itemize}





% Experimentación y Resultados
% \include{components/experimentacion/resultados}

%Referencias
\begin{thebibliography}{9}

  \bibitem{herrero2017fisica}
Herrero, M. and Escartín, J. C. G. (2017). 
\emph{La física de los números aleatorios}. 
Revista española de física, 31(2), 23--26.

\bibitem{Nansai2013}
Nansai, S., Elara, M. R., \& Iwase, M. (2013). 
Dynamic Analysis and Modeling of Jansen Mechanism. 
\textit{Procedia Engineering}, \textbf{64}, 1562--1571. 
Elsevier. 
https://doi.org/10.1016/J.PROENG.2013.09.238

  
\bibitem{PachecoHernndez2021}
Pacheco Hernández, P. R., Salini Calderón, G. A., \& Mera Garrido, E. M. (2021). 
\textit{Entropía y neguentropía: Una aproximación al proceso de difusión de contaminantes y su sostenibilidad}. 
Revista Internacional de Contaminación Ambiental, 37, 167--185. 
https://doi.org/10.20937/RICA.53688

\bibitem{herrera2000numeros}
Herrera, A. M. M. (2000). Números aleatorios: Historia, teoría y aplicaciones. \textit{Ingeniería y Desarrollo}, (8), 49--69.

\bibitem{seguridadinformatica}
\textit{Introducción a la seguridad informática y el análisis de vulnerabilidades}. Recuperado de \url{https://play.google.com/books/reader?id=5Z9yDwAAQBAJ&pg=GBS.PA16&hl=es_419}


\bibitem{tellez2016diseno}
Téllez, J. (2016). Diseño de mecanismo de locomoción andante con cambio de dirección. \textit{Navarra}, junio.

  \bibitem{OrdenOlasagasti}
  Olga del Orden Olasagasti. \textit{Ejercicios Resueltos De Inferencia Estadística Y Del Modelo Lineal Simple}. Recuperado de \url{https://books.google.com.pe/books?hl=es&lr=&id=pr5TEQAAQBAJ&oi=fnd&pg=PR2&dq=Ejercicios+resueltos+de+inferencia+estad%C3%ADstica+y+del+modelo+lineal+simple&ots=T42RpsBjJu&sig=jAkyfuPBzL1v0o3tHDlm0S5LbGI&redir_esc=y#v=onepage&q=Ejercicios%20resueltos%20de%20inferencia%20estad%C3%ADstica%20y%20del%20modelo%20lineal%20simple&f=false}.

  \bibitem{GilIzquierdo2014}
  Gil Izquierdo, M., Gonzáles Martín, A. I., \& Jano Salagre, M. D. (2014). \textit{Ejercicios de estadística teórica: probabilidad e Inferencia (2a ed.)}. Editorial Universidad Autónoma de Madrid.

  \bibitem{DiazRodriguez}
  Díaz Rodríguez, Martín. \textit{Estadística inferencial aplicada: Segunda edición revisada y aumentada}. Recuperado de \url{https://books.google.com.pe/books?hl=es&lr=&id=IHCHEAAAQBAJ&oi=fnd&pg=PP1&dq=Estad%C3%ADstica+inferencial+aplicada++Autor:+D%C3%ADaz+Rodr%C3%ADguez,+Mart%C3%ADn+-+Autor&ots=Sf8Tno6kLg&sig=u1xdyAJline4t90FJ5YAAxxWKng&redir_esc=y#v=onepage&q=Estad%C3%ADstica%20inferencial%20aplicada%20%20Autor%3A%20D%C3%ADaz%20Rodr%C3%ADguez%2C%20Mart%C3%ADn%20-%20Autor&f=false}.

  \bibitem{LlinasSolano}
  Llinás Solano, Humberto. \textit{Estadística inferencial}. Recuperado de \url{https://books.google.com.pe/books?hl=es&lr=&id=KZYQEAAAQBAJ&oi=fnd&pg=PA1&dq=Estad%C3%ADstica+Inferencial++Autor:+Llin%C3%A1s+Solano,+Humberto&ots=Xpi0wWw3lo&sig=DWM-pZnBeIOjGyyL5sBk_s_LWBQ&redir_esc=y#v=onepage&q=Estad%C3%ADstica%20Inferencial%20%20Autor%3A%20Llin%C3%A1s%20Solano%2C%20Humberto&f=false}.

  \bibitem{GomezVillegas}
  Gómez Villegas, Miguel Ángel. \textit{Inferencia estadística}. Recuperado de \url{https://books.google.com.pe/books?hl=es&lr=&id=YOuODwAAQBAJ&oi=fnd&pg=PR9&dq=Miguel+%C3%81ngel+G%C3%B3mez+Villegas+ESTAD%C3%8DSITICA+INFERENCIAL&ots=KG2bqmMDNd&sig=xguJOxmXHDguzSQsmn17Gy-SlNA&redir_esc=y#v=onepage&q=Miguel%20%C3%81ngel%20G%C3%B3mez%20Villegas%20ESTAD%C3%8DSITICA%20INFERENCIAL&f=false}.


\end{thebibliography}

\end{document}